\section{Experimental Protocol in Full}
\label{app:experiments}

This appendix gives the details a reader would need to re-run what we ran. It also
lists the exclusions and the frozen protocols by name, since several of our
reported failures are only interpretable against the criterion they failed.

\subsection{Stream}
\label{app:stream}

The Order-4 sequence of $15$ tasks, in presentation order: MNLI, CB, WiC, COPA,
QQP, BoolQA, RTE, IMDB, Yelp, Amazon, SST-2, DBpedia, AGNews, MultiRC, Yahoo. Each
task is rendered as a prompt containing an explicit option list, and the target is
a verbalizer token. The option list is readable at inference; the task index is
not, and no arm in this paper receives one except the oracle rung, which exists
only as an upper bound.

Scope structure follows directly from the verbalizer sets and is given in
Section~\ref{sec:setup}: four multi-task scopes and five singletons, with ten of
the fifteen tasks sharing output tokens with at least one other task.
Table~\ref{tab:scope-dims} in Appendix~\ref{app:proof-qm} records the gauge
dimension $d=K_S-1$ of each multi-task scope, which is what decides whether our
$q_m$ results are covered by proof.

\subsection{Model, budget, and optimisation}
\label{app:model}

T5-large with LoRA rank $8$ on the $q$ and $v$ projections of every attention
block: $2{,}359{,}296$ trainable scalars. The count is identical in every arm and
is asserted at run time per arm rather than assumed. The backbone and the
unembedding are frozen throughout; no arm trains an output head, and the
verbalizer rows $\bu_y$ are never updated.

Per task: $400$ examples per class, $3$ epochs, effective batch $64$. On the
$15$-task stream this is $813$ gradient steps per seed. Three seeds throughout,
and every reported level is a mean or median over seeds as stated at the point of
use.

\paragraph{A parameter that was dead, and the regression test that pins it.}
An earlier version of the harness passed the epoch count through a path that did
not reach the trainer, so the full $15$-task stream ran on $74$ gradient steps in
total. Every level measured under that configuration is uninterpretable as a
statement about a trained model, and the phase that depended on it failed its
second criterion for that reason. The parameter is fixed, and a regression test
asserts the realised step count against the configured one so the failure mode
cannot recur silently. We mention it here because it is the kind of defect that is
invisible in a results table.

\subsection{Splits}
\label{app:splits}

Each task's training file is partitioned into three disjoint sets:
\begin{itemize}
  \item \texttt{update} --- training, capped at $400$/class;
  \item \texttt{risk} --- $64$/class, used to \emph{fit} offsets;
  \item \texttt{audit} --- $64$/class, used to \emph{score} them.
\end{itemize}
Offsets are fit on \texttt{risk} and scored on \texttt{audit}, and \texttt{risk} is
never reported as a level. \textbf{The official test files are never opened at any
point in this paper.} The $64$/class audit split is balanced by construction, which
is a confound for any method whose behaviour depends on the label mix of the batch
it sees; where that applies we say so and test it directly rather than leaving it
to the reader.

\subsection{Exclusions, and why each one is not a choice we made after seeing scores}
\label{app:exclusions}

Three tasks are excluded from scored means. In each case the reason is a property
of the task's verbalizer or its class counts, checkable before any model is trained:

\begin{itemize}
  \item \textbf{Yelp} and \textbf{Amazon}: their five-way verbalizers violate the
    distinct-first-piece precondition. Under T5's sentencepiece tokenizer
    \texttt{very negative} and \texttt{very positive} share their first piece
    (id $182$), so a first-step restricted argmax cannot separate them, and the
    quantity we would be scoring is not the quantity we define.
  \item \textbf{CB}: its rarest class (\texttt{neutral}) has $16$ examples, below
    the $40$ that a $64$/class risk-plus-audit split requires. The frozen protocol
    specifies $64$/class, so CB cannot supply one. The cost of this exclusion is
    real and we state it: the NLI group reduces to MNLI and RTE.
\end{itemize}

Twelve of fifteen tasks remain scorable. Excluded tasks are still \emph{trained
on}, and they still compete for the shared offset --- omitting them from conflict
statistics would understate the conflict --- so they appear in $q_m$ and in the
scope tables while being absent from accuracy means.

\subsection{Metrics}
\label{app:metrics}

All accuracies are \textbf{balanced accuracy} under the restricted argmax of
\eqref{eq:raw-decode}, i.e.\ the mean of per-class recalls over the task's
verbalizer set. Balanced accuracy is used because the exclusion rules leave tasks
with unequal class counts and because a shared offset moves predicted class
proportions, which unweighted accuracy would partly hide.

The three retentions $R^{\mathrm{raw}}$, $R^{\mathrm{orc}}$, $R^{\mathrm{shr}}$ are
defined in \eqref{eq:orc}--\eqref{eq:shr}, the identification gap in
\eqref{eq:delta-id}, and the ladder rungs in Section~\ref{sec:ladder}. Offsets are
gauge-fixed to zero mean on their scope before being stored, applied, or compared,
per Appendix~\ref{app:gauge}.

\subsection{Statistics}
\label{app:stats}

All intervals are cluster bootstrap with \textbf{tasks as clusters}, $10{,}000$
draws. Tasks are the clusters because the same task measured at different stream
positions shares both its data and the model state that produced it; resampling
examples would treat those observations as independent and understate variance.
Where a comparison is paired by construction --- two inference rules evaluated on
the same task at the same stage --- the bootstrap resamples the paired differences.

Point estimates are medians unless stated otherwise. We report the interval
whenever we report the point, including for criteria that failed, and we do not
convert a failed criterion into a tie.

\subsection{The unconditional training gate}
\label{app:gate}

Every arm reports median post-training $R^{\mathrm{raw}}$, and the gate is
$>0.60$. The gate is unconditional: it is reported whatever its value, and it is
reported before the arm's forgetting numbers. Below the gate an arm has not learned
the stream and nothing about its forgetting is interpretable --- an arm that fails
the gate and shows little forgetting has shown nothing. Every arm in this paper
passes.

\subsection{Frozen protocols}
\label{app:protocols}

Each phase has a protocol file written before the run it governs, containing the
method, the criteria with numeric thresholds, and a set of predeclared outcomes
mapping possible results to conclusions. The file's SHA-256 is recorded on disk at
freeze time. Amendments are \emph{appended} with their own date and a new hash;
nothing is edited in place, so a criterion cannot be restated after its result is
known.

Two properties of this scheme are worth naming because they bound what it
guarantees. First, the hashes pin the artefact, and the ordering of freeze and run
rests on the committed source and the appended amendment record rather than on a
timestamped external registry; we say so rather than implying a stronger form of
pre-registration than we have. Second, the decision scripts are themselves frozen
before results land, with thresholds as module-level constants and a constructed
failure case per criterion in their test suites, so the mapping from numbers to
verdicts is fixed before the numbers exist.

\paragraph{Where the scheme left a gap, in the phase that tested our own predictions.}
The sensitivity analysis of Section~\ref{sec:sso} was frozen with six criteria and
five predeclared outcomes, and \emph{none of the five fired verbatim}: every outcome
that mentioned the staleness criterion presupposed it passing, and it failed. The
decision script emits a null outcome and reports each criterion under its own name,
which is what the protocol's general rule requires; we record the gap rather than
adding a sixth branch after the fact. Two further disclosures belong with it. The
handling of degenerate bootstrap resamples --- draws whose pooled regressor has no
variation, so the slope is undefined --- was \emph{not} in the frozen text; we drop
such draws and report their fraction (between $0.01\%$ and $0.20\%$ here) rather than
scoring them as zero slope, which would disguise ``undefined'' as ``no effect'', and
that choice post-dates seeing the first point estimates. And one arm showed a
significant scale dependence that we then traced to a mechanical artefact: at a
single-member scope the scoped gap is zero \emph{by Proposition~\ref{prop:disjoint}}
rather than by measurement, so including those points anchors any regression on
membership count at zero. Restricting to multi-member scopes removes the significance
in every arm. That diagnosis is post hoc, so we use it only to decline a finding,
never to rescue one.

\subsection{What was run on what}
\label{app:compute}

All runs are on a single A100-80GB per seed, three seeds in parallel on separate
cards of a shared host. The backbone is loaded in bf16 with gradient checkpointing,
which is what makes a T5-large LoRA run fit alongside other tenants; the practical
per-card headroom on the shared host was between $3.5$ and $6$~GB, and the
configuration above was sized to that rather than to an empty machine. A full
$15$-task seed takes roughly one hour, with later stages slower than earlier ones
because every stage re-evaluates every task seen so far.
